\uuid{EiCX}
\titre{Implication/équivalence}
\auteur{bodin ; hanani ; mzari}
\organisation{exo7}
\licence{CC BY-NC-SA 4.0}
\datecreate{2020-01-20}
\niveau{L1}
\module{Algèbre}
\chapitre{Logique, ensemble, raisonnement}
\sousChapitre{Absurde et contraposée}
\difficulte{3}
\theme{Raisonnement, implication, équivalence}
\exo7qcmid{lille-1.logique.027}

\contenu{
\begin{question}
\qclassification{100.03,100.04}
\qtags{motcle=implication/équivalence}
Le raisonnement par contraposée est basé
sur le fait que "$P \implies Q$" est équivalent à:
\begin{answers}
    \bad{"non($P$) $\implies$ non($Q$)".}

    \good{"non($Q$) $\implies$ non($P$)".}

    \bad{"non($P$) ou $Q$".}

    \bad{"$P$ ou non($Q$)".}
\end{answers}
\begin{explanations}
La contraposée de "$P \implies Q$" est "non($Q$) $\implies$ non($P$)".
\end{explanations}
\end{question}
}
